\section{Method}
\label{sec:method}

\subsection{IPDA Debate Format}

\begin{figure*}[t]
\centering
\includegraphics[width=\textwidth]{fig2_debate_sequence.pdf}
\caption{IPDA debate sequence: seven speeches alternating between
affirmative and negative with two cross-examination periods, each demanding
different capabilities (case generation, strategic response, dialogue).}
\label{fig:debate_sequence}
\end{figure*}

IPDA consists of 7 speeches alternating between affirmative and negative
(two constructives, two cross-examinations, three rebuttals;
Figure~\ref{fig:debate_sequence}). This structure is a natural curriculum:
constructives require case generation, rebuttals require strategic
response, and CX requires real-time dialogue, motivating our group-wise
training. Arguments follow typed syllogistic structures drawn from standard
debate theory~\cite{snider2008code}.

\subsection{Debate Generation Pipeline}
\label{sec:pipeline}

Before any training innovation can operate, the system must first produce
a complete, grounded debate from only a resolution. Our preparation
pipeline is the substrate on which branching, hierarchical reward, and
group-wise GRPO all operate, and it is shared identically between baseline
and experimental models, so downstream quality differences come from the
speech generation model, not from prep asymmetries.

\textbf{Belief tree.} A DSPy \textsc{ValueExtractor} produces 4
topic-specific values (2 per side); a \textsc{BeliefGenerator} expands
each to a depth-2 tree ($\sim$32 beliefs, $\sim$16 arguments per side).
Each belief carries a Bayesian credence in $[0,1]$, and arguments are
typed (uniqueness, link, impact, solvency, turn) for structured downstream
slots. \textbf{Research loop.} Per argument claim: LLM query generation
$\rightarrow$ parallel Weaviate$+$Tavily retrieval $\rightarrow$ sentence
chunking $\rightarrow$ hybrid scoring ($0.7\cdot\cos + 0.3\cdot$BM25)
$\rightarrow$ window expansion $\rightarrow$ structured LLM selection.
This consolidated loop uses \textbf{2 LLM calls per hop} vs.\ 9 in prior
systems, with most claims resolving in one hop. \textbf{Evidence cards}
store \texttt{full\_text}, \texttt{card} (with \textbf{bolded} selections),
\texttt{selected\_sentence\_ids}, and window ranges; speech generation
cites by sentence~ID only, so quotations are assembled verbatim and
fabrication is eliminated by construction. \textbf{Perspective building.}
A \textsc{PerspectiveBuilder} compiles side-specific prompts:
constructives get pruned depth-1 beliefs (2 of 4); rebuttals get the full
sub-belief tree in SEAM form (Setup, Evidence, Analysis, Move).
\textbf{Four-stage speech generation.} Each of the 7 speeches runs
(1)~\textsc{SelectTactic}, (2)~\textsc{BuildSkeleton},
(3)~\textsc{SelectEvidence} (constructives only; binds claims to
sentence~IDs), and (4)~\textsc{GenerateSpeech}. The training innovations
described next operate on calls produced by this pipeline. Full schemas
are in Appendix~\ref{app:pipeline}.

\subsection{Branching Debate Generation}
\label{sec:branching}

\begin{figure*}[t]
\centering
\includegraphics[width=\textwidth]{fig_branching_tree.pdf}
\caption{Branching debate tree. At each speech point, $N$ trials are
generated and scored; the debate forks when trials exceed quality
threshold and diverge enough that the opponent would respond differently.
Winner-shifts between branches yield counterfactual training pairs.}
\label{fig:branching_tree}
\end{figure*}

Standard debate generation is linear: produce one speech, continue to the
next. We instead treat debate generation as \textbf{search over debate
trajectories} (Figure~\ref{fig:branching_tree}). At each speech point, the model
generates $N = 3$ candidate trials scored by dimensional judges. When two
or more trials exceed a quality threshold \emph{and} diverge sufficiently
that the opponent would respond differently, the debate forks into parallel
paths (capped at 4 branches). The critical signal arises from
\textbf{winner-shift detection}: when branches from the same decision point
produce different debate outcomes, we obtain counterfactual training pairs
that receive $2\times$ weight in GRPO (Section~\ref{sec:grpo}), isolating
the specific decision that changed the result. Before each speech, the model
selects a named tactic via an epsilon-greedy policy, serving as a commitment
device that improves coherence. Full branching thresholds and the state
tracking algorithm are detailed in Appendix~\ref{app:branching}.

\subsection{Hierarchical Reward System}
\label{sec:reward_hierarchy}

\begin{figure*}[t]
\centering
\includegraphics[width=\textwidth]{fig_reward_hierarchy.pdf}
\caption{Three-layer hierarchical reward. Layer~1: dimensional speech
scoring. Layer~2: outcome backpropagation to calls. Layer~3: multi-model
retrospective analysis for bias and turning points.}
\label{fig:reward_hierarchy}
\end{figure*}

Scoring isolated speeches is insufficient: a sound rebuttal can still lose
if it addresses the wrong arguments. We employ a three-layer reward
hierarchy (Figure~\ref{fig:reward_hierarchy}).

\paragraph{Layer 1: Dimensional scoring.}
Four sub-judges (ArgumentJudge, EvidenceJudge, ClashJudge, LanguageJudge)
score each speech with specialized rubrics, plus a speech-specific
dimension reflecting strategic role (e.g., 1AR \emph{triage efficiency},
2AR \emph{crystallization}; Appendix~\ref{app:agentic_judge}).

\paragraph{Layer 2: Outcome backpropagation.}
A BackpropScorer propagates debate outcomes to individual calls with
causal attribution, assigning each speech a contribution in $[-1.0, +1.0]$.
Winner-shift events from branching produce direct adjustments at the
divergence point, so a strong speech in a losing debate is not punished
if the loss originated earlier.

\paragraph{Layer 3: Meta-debate comparison.}
A 3-model retrospective panel evaluates complete traces, detecting bias
and identifying turning points that become highest-signal training data.

\paragraph{Evidence verification.}
Models cite by sentence ID, assembling verbatim text from research
documents---eliminating fabrication by construction. Evidence quality
uses a 6-level scale with $2\times$ asymmetric GRPO penalty on
fabrications.

\subsection{Group-Wise Offline GRPO}
\label{sec:grpo}

\begin{algorithm}[t]
\caption{Group-Wise Offline GRPO for Debate}
\label{alg:training}
\begin{algorithmic}[1]
\FOR{$i = 1$ to $N$ iterations}
    \STATE Generate debates with branching; capture log-probs
    \STATE Score with 3-layer reward hierarchy
    \FOR{group $\in \{A, B, C, D\}$}
        \STATE Generate Opus golden samples for group
        \STATE SFT on golden $\rightarrow$ merge to SFT$^+$
        \STATE Generate new trials with SFT$^+$; capture log-probs
        \STATE Offline GRPO with importance sampling
        \STATE Merge LoRA $\rightarrow$ canonical model
    \ENDFOR
\ENDFOR
\end{algorithmic}
\end{algorithm}

A debate produces $\sim$33 LLM calls spanning radically different output
modalities: $\sim$200-token tactic selections, $\sim$2{,}000-token
speeches, and multi-turn CX. Training jointly creates noisy gradients as
the optimizer tries to satisfy contradictory length and format
distributions. We \textbf{decompose calls into four groups} processed
sequentially each iteration (Algorithm~\ref{alg:training}):

\begin{table}[t]
\centering
\small
\begin{tabular}{llcc}
\toprule
\textbf{Group} & \textbf{Call Types} & \textbf{Samples} & \textbf{Output Length} \\
\midrule
A & \textsc{tactic\_select} & $\sim$3K & Short \\
B & \textsc{skeleton\_build}, \textsc{evidence\_select} & $\sim$4K & Medium \\
C & \textsc{speech\_generate} & $\sim$4.5K & Long \\
D & \textsc{cx\_dialogue} & $\sim$5K & Dialogue \\
\bottomrule
\end{tabular}
\caption{Call-type groups for sequential GRPO training. Each group contains
calls with similar output modality and length distribution, preventing
cross-modality gradient interference.}
\label{tab:call_types}
\end{table}

For each group, we: (1)~generate Opus golden samples as quality upper
bounds, (2)~SFT on golden samples via LoRA ($r = 32$) to produce an SFT$^+$
checkpoint, (3)~generate new trials with SFT$^+$ while capturing
log-probabilities for the behavior policy, (4)~run offline
GRPO~\cite{shao2024deepseekmath} with importance sampling, and (5)~merge
the LoRA into the canonical model as baseline for the next group.
Advantages are normalized within \emph{(our\_tactic, opponent\_tactic,
opponent\_intent)} triplets rather than globally, yielding cleaner gradient
signal. Full details are in Appendix~\ref{app:grpo_groups}.

\paragraph{Judge design.}
Beyond the four Layer~1 sub-judges, a \textbf{Chronicler} tracks each
argument's status (standing, attacked, dropped, extended, conceded, turned)
across the debate, and a \textbf{Forensics module} attributes failures to
specific pipeline stages (tactic, skeleton, evidence, execution). For each
call, $N = 3$ trials at varied temperatures produce chosen/rejected
preference pairs, yielding $\sim$60+ pairs per debate
(Appendix~\ref{app:agentic_judge}).
